\chapter{THIẾT KẾ HỆ THỐNG}

\section{Sơ đồ lớp (Class Diagram)}

Dựa trên các thực thể nghiệp vụ đã phân tích, hệ thống được thiết kế với các lớp chính sau:

\begin{itemize}[leftmargin=1.5cm]
    \item \textbf{NguoiDung}: lưu thông tin tài khoản, quyền hạn.
    \item \textbf{HangHoa}: lưu thông tin hàng hóa, nhóm hàng, đơn vị tính.
    \item \textbf{NhaCungCap}: lưu thông tin nhà cung cấp.
    \item \textbf{KhachHang}: lưu thông tin khách hàng.
    \item \textbf{Kho}: lưu thông tin kho hàng.
    \item \textbf{PhieuNhapKho}: lưu thông tin phiếu nhập kho.
    \item \textbf{ChiTietPhieuNhap}: lưu chi tiết hàng hóa trong phiếu nhập.
    \item \textbf{PhieuXuatKho}: lưu thông tin phiếu xuất kho.
    \item \textbf{ChiTietPhieuXuat}: lưu chi tiết hàng hóa trong phiếu xuất.
    \item \textbf{PhieuKiemKe}: lưu thông tin phiếu kiểm kê.
    \item \textbf{ChiTietKiemKe}: lưu chi tiết kiểm kê từng hàng hóa.
    \item \textbf{TonKho}: lưu số lượng tồn kho của từng hàng hóa tại từng kho.
\end{itemize}

\begin{figure}[H]
\centering
\begin{tikzpicture}[scale=0.8, transform shape]

\begin{class}[text width=2.8cm]{NhomHangHoa}{-6,3}
\attribute{+ MaNhomHang: String}
\attribute{+ TenNhomHang: String}
\attribute{+ MoTa: String}
\end{class}

\begin{class}[text width=3.2cm]{HangHoa}{-1,3}
\attribute{+ MaHangHoa: String}
\attribute{+ TenHangHoa: String}
\attribute{+ MaNhomHang: String}
\attribute{+ MaDonViTinh: String}
\attribute{+ GiaNhap: Decimal}
\attribute{+ GiaXuat: Decimal}
\attribute{+ TrangThai: Boolean}
\operation{+ CapNhatThongTin(): void}
\end{class}

\begin{class}[text width=2.8cm]{DonViTinh}{4,3}
\attribute{+ MaDonViTinh: String}
\attribute{+ TenDonViTinh: String}
\end{class}

\begin{class}[text width=2.8cm]{Kho}{-6,-1}
\attribute{+ MaKho: String}
\attribute{+ TenKho: String}
\attribute{+ DiaChi: String}
\attribute{+ NguoiPhuTrach: String}
\attribute{+ DienTich: Decimal}
\attribute{+ TrangThai: Boolean}
\end{class}

\begin{class}[text width=3.2cm]{NguoiDung}{-1,-1}
\attribute{+ MaNguoiDung: String}
\attribute{+ TenDangNhap: String}
\attribute{+ MatKhau: String}
\attribute{+ HoTen: String}
\attribute{+ VaiTro: String}
\attribute{+ TrangThai: Boolean}
\operation{+ DangNhap(): Boolean}
\operation{+ DangXuat(): void}
\end{class}

\begin{class}[text width=3.2cm]{NhaCungCap}{4,-1}
\attribute{+ MaNhaCungCap: String}
\attribute{+ TenNhaCungCap: String}
\attribute{+ DiaChi: String}
\attribute{+ SoDienThoai: String}
\attribute{+ Email: String}
\attribute{+ NguoiLienHe: String}
\end{class}

\begin{class}[text width=3.2cm]{KhachHang}{-6,-5}
\attribute{+ MaKhachHang: String}
\attribute{+ TenKhachHang: String}
\attribute{+ DiaChi: String}
\attribute{+ SoDienThoai: String}
\attribute{+ Email: String}
\attribute{+ MaSoThue: String}
\attribute{+ NguoiLienHe: String}
\end{class}

\begin{class}[text width=2.8cm]{TonKho}{-1,-5}
\attribute{+ MaKho: String}
\attribute{+ MaHangHoa: String}
\attribute{+ SoLuongTon: Integer}
\attribute{+ NgayCapNhat: DateTime}
\operation{+ CapNhatTonKho(): void}
\end{class}

% Quan hệ
\draw (NhomHangHoa) -- (HangHoa) node[midway, above, font=\tiny] {1..*};
\draw (DonViTinh) -- (HangHoa) node[midway, above, font=\tiny] {1..*};
\draw (HangHoa) -- (TonKho) node[midway, right, font=\tiny] {1..*};
\draw (Kho) -- (TonKho) node[midway, above, font=\tiny] {1..*};

\end{tikzpicture}
\caption{Sơ đồ lớp danh mục (Master Data)}
\label{fig:class-diagram-master}
\end{figure}


\input{diagrams/classdiagram_nhapxuat}

\input{diagrams/classdiagram_kiemke_donhang}


\section{Sơ đồ trình tự (Sequence Diagram)}

\subsection{Sequence diagram ``Nhập kho''}

\input{diagrams/sequence_nhapkho}

\subsection{Sequence diagram ``Xuất kho''}

\input{diagrams/sequence_xuatkho}

\section{Thiết kế cơ sở dữ liệu}

\subsection{Mô hình quan hệ tổng thể}

Cơ sở dữ liệu của hệ thống bao gồm các bảng chính sau:

\begin{longtable}{|p{4cm}|p{9cm}|}
\hline
\textbf{Bảng} & \textbf{Mô tả} \\
\hline
\endfirsthead
\hline
\textbf{Bảng} & \textbf{Mô tả} \\
\hline
\endhead
NguoiDung & Lưu thông tin tài khoản người dùng: mã, tên đăng nhập, mật khẩu, vai trò, trạng thái. \\
\hline
HangHoa & Lưu thông tin hàng hóa: mã, tên, mã nhóm, đơn vị tính, giá nhập, giá xuất, mô tả. \\
\hline
NhomHangHoa & Lưu thông tin nhóm hàng hóa: mã, tên, mô tả. \\
\hline
DonViTinh & Lưu đơn vị tính: mã, tên. \\
\hline
NhaCungCap & Lưu thông tin nhà cung cấp: mã, tên, địa chỉ, số điện thoại, email. \\
\hline
KhachHang & Lưu thông tin khách hàng: mã, tên, địa chỉ, số điện thoại, email. \\
\hline
Kho & Lưu thông tin kho: mã, tên, địa chỉ, người phụ trách. \\
\hline
PhieuNhapKho & Lưu phiếu nhập kho: mã, mã nhà cung cấp, mã kho, ngày nhập, người lập, trạng thái, tổng tiền. \\
\hline
ChiTietPhieuNhap & Lưu chi tiết nhập: mã phiếu nhập, mã hàng hóa, số lượng, đơn giá, thành tiền. \\
\hline
PhieuXuatKho & Lưu phiếu xuất kho: mã, mã khách hàng, mã kho, ngày xuất, người lập, trạng thái, tổng tiền. \\
\hline
ChiTietPhieuXuat & Lưu chi tiết xuất: mã phiếu xuất, mã hàng hóa, số lượng, đơn giá, thành tiền. \\
\hline
PhieuKiemKe & Lưu phiếu kiểm kê: mã, mã kho, ngày kiểm, người kiểm, trạng thái. \\
\hline
ChiTietKiemKe & Lưu chi tiết kiểm kê: mã phiếu kiểm, mã hàng hóa, số lượng thực tế, số lượng sổ sách, chênh lệch. \\
\hline
TonKho & Lưu tồn kho: mã kho, mã hàng hóa, số lượng tồn, ngày cập nhật. \\
\hline
LichSuThaoTac & Lưu log thao tác: mã, người dùng, hành động, thời gian, đối tượng. \\
\hline
\end{longtable}

\subsection{Mô tả một số bảng quan trọng}

\subsubsection{Bảng HangHoa}

\begin{longtable}{|p{3.5cm}|p{3cm}|p{5.5cm}|}
\hline
\textbf{Thuộc tính} & \textbf{Kiểu dữ liệu} & \textbf{Mô tả} \\
\hline
\endfirsthead
\hline
\textbf{Thuộc tính} & \textbf{Kiểu dữ liệu} & \textbf{Mô tả} \\
\hline
\endhead
MaHangHoa & VARCHAR(20) & Khóa chính, mã hàng hóa. \\
\hline
TenHangHoa & NVARCHAR(255) & Tên hàng hóa. \\
\hline
MaNhomHang & VARCHAR(20) & Khóa ngoại đến bảng NhomHangHoa. \\
\hline
MaDonViTinh & VARCHAR(20) & Khóa ngoại đến bảng DonViTinh. \\
\hline
GiaNhap & DECIMAL(18,2) & Giá nhập trung bình. \\
\hline
GiaXuat & DECIMAL(18,2) & Giá xuất bán. \\
\hline
MoTa & NVARCHAR(500) & Mô tả chi tiết. \\
\hline
TrangThai & BIT & Trạng thái còn sử dụng. \\
\hline
\end{longtable}

\subsubsection{Bảng TonKho}

\begin{longtable}{|p{3.5cm}|p{3cm}|p{5.5cm}|}
\hline
\textbf{Thuộc tính} & \textbf{Kiểu dữ liệu} & \textbf{Mô tả} \\
\hline
\endfirsthead
\hline
\textbf{Thuộc tính} & \textbf{Kiểu dữ liệu} & \textbf{Mô tả} \\
\hline
\endhead
MaKho & VARCHAR(20) & Khóa ngoại đến bảng Kho. \\
\hline
MaHangHoa & VARCHAR(20) & Khóa ngoại đến bảng HangHoa. \\
\hline
SoLuongTon & INT & Số lượng tồn kho hiện tại. \\
\hline
NgayCapNhat & DATETIME & Thời gian cập nhật gần nhất. \\
\hline
\end{longtable}

\section{Kết luận chương}

Chương 3 đã trình bày thiết kế hệ thống ở mức độ chi tiết hơn, bao gồm sơ đồ lớp, sơ đồ trình tự cho các nghiệp vụ nhập kho và xuất kho, cùng với thiết kế cơ sở dữ liệu. Các thiết kế này là cơ sở để triển khai phần mềm quản lý kho hàng trong giai đoạn tiếp theo.
